\chapter{The Ricci flow on surfaces}

\section*{1. Evolution equations on surfaces}

\begin{theorem}[Conformal change of scalar curvature]
\label{ch05-thm-conformal-scalar-curvature}
\source{chow-knopf-ricci-flow:page-0120}
\dcref{Theorem 5.3}
\group{evolution-equations-on-surfaces}

If \(g\) and \(h\) are metrics on a surface \(\M^2\) conformally related by
\[
g=e^{u}h,
\]
then their scalar curvatures satisfy
\[
R_g=e^{-u}\left(R_h-\Delta_h u\right).
\]
\end{theorem}

\begin{proof}
\source{chow-knopf-ricci-flow:page-0120}
\sketch This is the standard conformal curvature transformation formula in dimension two.
It follows by writing the Levi-Civita connection of \(g\) in terms of that of \(h\), then contracting the curvature tensor. Since the conformal factor is scalar, the formula simplifies to the stated expression.
\end{proof}

\begin{lemma}[Evolution of the Laplacian]
\label{ch05-lem-laplacian-evolution}
\source{chow-knopf-ricci-flow:page-0121}
\dcref{Lemma 5.4}
\group{evolution-equations-on-surfaces}

Let \(g(t)\) be a smooth \(1\)-parameter family of metrics on \(\mathcal M^n\). Then for any smooth function \(u\),
\[
\frac{\partial}{\partial t}(\Delta u)
=\Delta\!\left(\frac{\partial u}{\partial t}\right)
-\left\langle \frac{\partial g}{\partial t},\Hess u\right\rangle
\]
and, in particular, if \(g(t)\) satisfies the normalized Ricci flow on a surface, then
\[
\frac{\partial}{\partial t}\Delta=(R-r)\Delta.
\]
\end{lemma}

\begin{proof}
\source{chow-knopf-ricci-flow:page-0121}
\sketch Differentiate the local coordinate formula for \(\Delta\) and use the evolution of the inverse metric and Christoffel symbols. On a surface under the normalized Ricci flow, \(\partial_t g=(r-R)g\), so the tensorial correction reduces to multiplication by \(R-r\).
\end{proof}

\begin{corollary}[Evolution of the area form]
\label{ch05-cor-area-form-evolution}
\source{chow-knopf-ricci-flow:page-0122}
\dcref{Corollary 5.5}
\group{evolution-equations-on-surfaces}

If \(g(t)\) is a smooth \(1\)-parameter family of metrics on \(\mathcal{M}^n\) such that
\[
\frac{\partial}{\partial t}g = h,
\]
then the volume form evolves by
\[
\frac{\partial}{\partial t}d\mu = \frac12 \operatorname{tr}_g(h)\,d\mu.
\]
\end{corollary}

\begin{corollary}[Area preservation under normalized flow]
\label{ch05-cor-area-preservation}
\source{chow-knopf-ricci-flow:page-0122}
\dcref{Corollary 5.6}
\group{evolution-equations-on-surfaces}

If \((\mathcal{M}^n,g(t))\) is a solution of the normalized Ricci flow
\[
\frac{\partial}{\partial t}g=(r-R)g,
\]
then the total volume \(\Vol(\mathcal M^n,g(t))\) is constant in time.
\end{corollary}

\begin{lemma}[Scalar curvature evolution on a surface]
\label{ch05-lem-scalar-curvature-evolution-surface}
\source{chow-knopf-ricci-flow:page-0122}
\dcref{Lemma 5.7}
\group{evolution-equations-on-surfaces}

If \(g(t)\) is a smooth \(1\)-parameter family of metrics on a Riemannian surface \(\M^2\) such that \(\partial_t g = f g\) for a scalar function \(f\), then
\[
\frac{\partial}{\partial t}R = -fR -\Delta f.
\]
In particular, under the normalized Ricci flow,
\[
\frac{\partial R}{\partial t}=\Delta R+R(R-r).
\]
\end{lemma}

\begin{proof}
\source{chow-knopf-ricci-flow:page-0122, chow-knopf-ricci-flow:page-0123}
\sketch The first formula is the two-dimensional specialization of the general scalar curvature variation formula. Substituting \(f=r-R\) gives the normalized Ricci flow evolution equation.
\end{proof}

\begin{corollary}[Evolution of the curvature potential]
\label{ch05-cor-curvature-potential-evolution}
\source{chow-knopf-ricci-flow:page-0123}
\dcref{Corollary 5.8}
\group{evolution-equations-on-surfaces}

If \(g(t)\) is a solution of the normalized Ricci flow on a surface and \(f\) is the curvature potential defined by
\[
\Delta f = R-r,
\]
then
\[
\frac{\partial f}{\partial t}=\Delta f + r f - \frac12 |\nabla f|^2 + c(t)
\]
for a function of time \(c(t)\).
\end{corollary}

\begin{lemma}[Potential normalization]
\label{ch05-lem-potential-normalization}
\source{chow-knopf-ricci-flow:page-0124}
\dcref{Lemma 5.9}
\group{evolution-equations-on-surfaces}

Let \(g(t)\) be a complete solution with bounded curvature of the normalized Ricci flow on a surface. Then the potential \(f\) may be normalized uniquely by requiring
\[
\int_{\M^2} f\,d\mu = 0.
\]
\end{lemma}

\section*{2. The normalized Ricci flow}

\begin{definition}[Normalized flow]
\label{ch05-def-normalized-flow}
\source{chow-knopf-ricci-flow:page-0125}
\dcref{Definition 5.10}
\group{normalized-ricci-flow}

A solution \(g(t)\) of the normalized Ricci flow on a surface is a smooth family of metrics satisfying
\[
\frac{\partial}{\partial t}g=(r-R)g.
\]
\end{definition}

\begin{definition}[Ricci potential]
\label{ch05-def-ricci-potential}
\source{chow-knopf-ricci-flow:page-0125}
\dcref{Definition 5.11}
\group{normalized-ricci-flow}

A metric \(g(t)\) on a surface \(\M^2\) is said to admit a Ricci potential \(f\) if
\[
\Delta f=R-r.
\]
\end{definition}

\begin{lemma}[Potential equation]
\label{ch05-lem-potential-equation}
\source{chow-knopf-ricci-flow:page-0126}
\dcref{Lemma 5.12}
\group{normalized-ricci-flow}

For a solution \(g(t)\) of the normalized Ricci flow on a compact surface with curvature potential \(f\),
\[
\frac{\partial f}{\partial t}=\Delta f + r f.
\]
\end{lemma}

\begin{remark}
\label{ch05-rem-potential-normalization}
\source{chow-knopf-ricci-flow:page-0126}
\dcref{Remark 5.13}
\group{normalized-ricci-flow}

This normalization for \(f\) differs from the one used in [60].
\end{remark}

\begin{corollary}[Time evolution of the potential]
\label{ch05-cor-potential-time-evolution}
\source{chow-knopf-ricci-flow:page-0127}
\dcref{Corollary 5.14}
\group{normalized-ricci-flow}

Under the normalized Ricci flow on a compact surface, there exists a constant \(C\) such that the potential \(f\) satisfies uniform integral bounds depending only on the initial metric.
\end{corollary}

\begin{proposition}[Curvature bounds for \(r\le 0\)]
\label{ch05-prop-curvature-bounds-nonpositive-average}
\source{chow-knopf-ricci-flow:page-0127}
\dcref{Proposition 5.15}
\group{normalized-ricci-flow}

Let \((\M^2,g(t))\) be a solution of the normalized Ricci flow on a compact surface with \(r\le 0\). Then there exists a constant \(C\ge 1\), depending only on \(g(0)\), such that for as long as the solution exists,
\[
\frac1C \le R(x,t)\le C
\]
and the metrics \(g(t)\) remain uniformly equivalent to \(g(0)\).
\end{proposition}

\begin{proposition}[Evolution of \(H\)]
\label{ch05-prop-H-evolution}
\source{chow-knopf-ricci-flow:page-0127}
\dcref{Proposition 5.16}
\group{normalized-ricci-flow}

On a solution \((\M^2,g(t))\) of the normalized Ricci flow on a compact surface, the quantity
\[
H:=R-r+\Delta f
\]
evolves by a scalar parabolic inequality that implies a time-dependent upper bound for \(R\).
\end{proposition}

\begin{corollary}[Upper curvature control]
\label{ch05-cor-upper-curvature-control}
\source{chow-knopf-ricci-flow:page-0128}
\dcref{Corollary 5.17}
\group{normalized-ricci-flow}

On a solution of the normalized Ricci flow on a compact surface, there exists a constant \(C\) depending only on the initial metric such that
\[
R(x,t)\le C
\]
for all \(x\) and all \(t\) in the existence interval.
\end{corollary}

\begin{proposition}[Two-sided curvature control]
\label{ch05-prop-two-sided-curvature-control}
\source{chow-knopf-ricci-flow:page-0128}
\dcref{Proposition 5.18}
\group{normalized-ricci-flow}

For any solution \((M^2,g(t))\) of the normalized Ricci flow on a compact surface, there exists a constant \(C>0\) depending only on the initial metric such that the scalar curvature is bounded above and below by constants depending only on \(C\) and \(t\).
\end{proposition}

\begin{proposition}[Long-time existence for \(r\le 0\)]
\label{ch05-prop-long-time-existence-nonpositive}
\source{chow-knopf-ricci-flow:page-0129}
\dcref{Proposition 5.19}
\group{normalized-ricci-flow}

If \((\M^2,g_0)\) is a closed Riemannian surface with \(r\le 0\), then the unique solution \(g(t)\) of the normalized Ricci flow with \(g(0)=g_0\) exists for all time.
\end{proposition}

\begin{proposition}[Self-similar solutions]
\label{ch05-prop-self-similar-solutions-constant-curvature}
\source{chow-knopf-ricci-flow:page-0130}
\dcref{Proposition 5.20}
\group{normalized-ricci-flow}

Any expanding or steady self-similar solution of the normalized Ricci flow on a closed surface has nonpositive average scalar curvature. In particular, if the solution is self-similar, then the metric is forced into the constant-curvature regime.
\end{proposition}

\begin{proposition}[Self-similarity implies constant curvature]
\label{ch05-prop-self-similar-constant-curvature}
\source{chow-knopf-ricci-flow:page-0131}
\dcref{Proposition 5.21}
\group{normalized-ricci-flow}

If \((\mathcal{M}^2,g(t))\) is a self-similar solution of the normalized Ricci flow on a Riemannian surface, then \(g(t)\equiv g(0)\) is a metric of constant curvature.
\end{proposition}

\section*{3. The cases \(r<0\) and \(r=0\)}

\begin{theorem}[Negative average curvature]
\label{ch05-thm-negative-average-curvature}
\source{chow-knopf-ricci-flow:page-0133, chow-knopf-ricci-flow:page-0134, chow-knopf-ricci-flow:page-0135, chow-knopf-ricci-flow:page-0136}
\dcref{Theorem 5.22}
\group{negative-and-zero-average-curvature}

If \((\M^2,g_0)\) is a closed Riemannian surface with average scalar curvature \(r<0\), then the unique solution \(g(t)\) of the normalized Ricci flow with \(g(0)=g_0\) converges uniformly in any \(C^k\)-norm to a smooth constant-curvature metric \(g_\infty\) as \(t\to\infty\).
\end{theorem}

\begin{theorem}[Zero average curvature]
\label{ch05-thm-zero-average-curvature}
\source{chow-knopf-ricci-flow:page-0136, chow-knopf-ricci-flow:page-0137, chow-knopf-ricci-flow:page-0138, chow-knopf-ricci-flow:page-0139, chow-knopf-ricci-flow:page-0140, chow-knopf-ricci-flow:page-0141}
\dcref{Theorem 5.28}
\group{negative-and-zero-average-curvature}

Let \((\mathcal{M}^2,g_0)\) be a closed Riemannian surface with average scalar curvature \(r=0\). Then the unique solution \(g(t)\) of the normalized Ricci flow with \(g(0)=g_0\) converges uniformly in any \(C^k\)-norm to a smooth constant-curvature metric \(g_\infty\) as \(t\to\infty\).
\end{theorem}

\begin{proposition}[Entropy monotonicity]
\label{ch05-prop-entropy-monotonicity}
\source{chow-knopf-ricci-flow:page-0142, chow-knopf-ricci-flow:page-0143}
\dcref{Proposition 5.39}
\group{positive-average-curvature}

Along a solution of the normalized Ricci flow on a surface, the entropy functional \(N(g(t))\) is nondecreasing, and it is strictly monotone unless the metric is Einstein.
\end{proposition}

\begin{proposition}[Positive curvature strategy]
\label{ch05-prop-positive-curvature-strategy}
\source{chow-knopf-ricci-flow:page-0144, chow-knopf-ricci-flow:page-0145}
\dcref{Proposition 5.44}
\group{positive-average-curvature}

For solutions on a surface with \(r>0\), the entropy and differential Harnack estimates combine with injectivity-radius control to yield uniform curvature bounds once the scalar curvature becomes positive.
\end{proposition}

\begin{lemma}[Doubling-time estimate]
\label{ch05-lem-doubling-time}
\source{chow-knopf-ricci-flow:page-0151}
\dcref{Lemma 5.45}
\group{positive-average-curvature}

Let \((\mathcal{M}^2,g(t))\) be any solution of the normalized Ricci flow. Then the scalar curvature cannot double too quickly: there is a universal lower bound on the time needed for \(\max R\) to increase by a factor of \(2\).
\end{lemma}

\begin{corollary}[Short-time curvature control]
\label{ch05-cor-short-time-curvature-control}
\source{chow-knopf-ricci-flow:page-0151}
\dcref{Corollary 5.47}
\group{positive-average-curvature}

There is a short-time upper bound for the scalar curvature of any normalized Ricci flow solution on a compact surface, expressed in terms of the initial maximum curvature.
\end{corollary}

\begin{lemma}[Gradient control]
\label{ch05-lem-gradient-control}
\source{chow-knopf-ricci-flow:page-0152}
\dcref{Lemma 5.48}
\group{positive-average-curvature}

For any solution of the normalized Ricci flow on a compact surface, the gradient of the scalar curvature satisfies a time-dependent estimate controlled by the initial metric.
\end{lemma}

\begin{corollary}[Curvature decay away from positivity]
\label{ch05-cor-curvature-decay-away-from-positivity}
\source{chow-knopf-ricci-flow:page-0153}
\dcref{Corollary 5.49}
\group{positive-average-curvature}

If the normalized Ricci flow on a compact surface is not immediately positively curved, then the scalar curvature still obeys a uniform estimate that depends only on the initial metric.
\end{corollary}

\begin{proposition}[Uniform curvature bound in the positive case]
\label{ch05-prop-uniform-curvature-bound-positive}
\source{chow-knopf-ricci-flow:page-0153}
\dcref{Proposition 5.50}
\group{positive-average-curvature}

There exists a universal constant \(C<\infty\) such that on any compact surface with \(r>0\), the normalized Ricci flow satisfies a uniform upper bound for the scalar curvature depending only on the initial metric.
\end{proposition}

\begin{proposition}[Diameter bound]
\label{ch05-prop-diameter-bound-positive}
\source{chow-knopf-ricci-flow:page-0154, chow-knopf-ricci-flow:page-0155}
\dcref{Proposition 5.51}
\group{positive-average-curvature}

For a solution of the normalized Ricci flow on a closed surface with \(R(\cdot,0)>0\), the diameter remains uniformly bounded for all time.
\end{proposition}

\begin{corollary}[Injectivity-radius control]
\label{ch05-cor-injectivity-radius-control}
\source{chow-knopf-ricci-flow:page-0156}
\dcref{Corollary 5.52}
\group{positive-average-curvature}

For any solution of the normalized Ricci flow on a closed surface with \(R(\cdot,0)>0\), there exists a constant \(C>0\) depending only on \(g_0\) such that the injectivity radius is bounded below uniformly in time.
\end{corollary}

\begin{corollary}[Differential Harnack inequality]
\label{ch05-cor-differential-harnack}
\source{chow-knopf-ricci-flow:page-0158}
\dcref{Corollary 5.56}
\group{positive-average-curvature}

The differential Harnack quantity \(Q\) satisfies a parabolic inequality that yields a lower bound depending only on the initial metric.
\end{corollary}

\begin{proposition}[Positive curvature lower bound]
\label{ch05-prop-positive-curvature-lower-bound}
\source{chow-knopf-ricci-flow:page-0158}
\dcref{Proposition 5.57}
\group{positive-average-curvature}

On a complete solution of the normalized Ricci flow with bounded positive scalar curvature, there exists a constant \(C>1\) depending only on \(g_0\) such that the scalar curvature is bounded below in terms of \(C\).
\end{proposition}

\begin{proposition}[Modified Harnack control]
\label{ch05-prop-modified-harnack-control}
\source{chow-knopf-ricci-flow:page-0159}
\dcref{Proposition 5.58}
\group{positive-average-curvature}

When the scalar curvature changes sign, an adjusted Harnack quantity remains controlled and can still be used to derive a positive lower bound after a finite time.
\end{proposition}

\begin{lemma}[Auxiliary lower-bound estimate]
\label{ch05-lem-aux-lower-bound}
\source{chow-knopf-ricci-flow:page-0159, chow-knopf-ricci-flow:page-0160}
\dcref{Lemma 5.59}
\group{positive-average-curvature}

The auxiliary quantity introduced for the sign-changing curvature case satisfies an evolution inequality that supports the lower-bound argument for \(R-s\).
\end{lemma}

\begin{proposition}[Entropy lower bound in the sign-changing case]
\label{ch05-prop-entropy-lower-bound-sign-changing}
\source{chow-knopf-ricci-flow:page-0161}
\dcref{Proposition 5.61}
\group{positive-average-curvature}

In the case where the initial scalar curvature is not everywhere positive, the entropy method still gives a uniform lower bound for the scalar curvature once the modified flow is used.
\end{proposition}

\begin{proposition}[Eventual positivity]
\label{ch05-prop-eventual-positivity}
\source{chow-knopf-ricci-flow:page-0161}
\dcref{Proposition 5.62}
\group{positive-average-curvature}

For a solution of the normalized Ricci flow on a surface of positive Euler characteristic, the scalar curvature eventually becomes strictly positive.
\end{proposition}

\begin{corollary}[Exponential convergence for \(r>0\)]
\label{ch05-cor-exponential-convergence-positive}
\source{chow-knopf-ricci-flow:page-0162}
\dcref{Corollary 5.63}
\group{positive-average-curvature}

If \((\M^2,g(t))\) is a solution of the modified Ricci flow associated to a surface with \(R(\cdot,0)>0\), then all curvature derivatives decay exponentially and the flow converges smoothly to a constant-curvature metric.
\end{corollary}

\begin{theorem}[Positive average curvature]
\label{ch05-thm-positive-average-curvature}
\source{chow-knopf-ricci-flow:page-0162, chow-knopf-ricci-flow:page-0163, chow-knopf-ricci-flow:page-0164, chow-knopf-ricci-flow:page-0165, chow-knopf-ricci-flow:page-0166, chow-knopf-ricci-flow:page-0167, chow-knopf-ricci-flow:page-0168, chow-knopf-ricci-flow:page-0169, chow-knopf-ricci-flow:page-0170, chow-knopf-ricci-flow:page-0171, chow-knopf-ricci-flow:page-0172, chow-knopf-ricci-flow:page-0173, chow-knopf-ricci-flow:page-0174, chow-knopf-ricci-flow:page-0175, chow-knopf-ricci-flow:page-0176, chow-knopf-ricci-flow:page-0177, chow-knopf-ricci-flow:page-0178, chow-knopf-ricci-flow:page-0179, chow-knopf-ricci-flow:page-0180, chow-knopf-ricci-flow:page-0181, chow-knopf-ricci-flow:page-0182, chow-knopf-ricci-flow:page-0183, chow-knopf-ricci-flow:page-0184, chow-knopf-ricci-flow:page-0185}
\dcref{Theorem 5.64}
\group{positive-average-curvature}

Let \((\M^2,g_0)\) be a closed Riemannian surface with average scalar curvature \(r>0\). The unique solution \(g(t)\) of the normalized Ricci flow with \(g(0)=g_0\) converges exponentially in any \(C^k\)-norm to a smooth constant-curvature metric \(g_\infty\) as \(t\to\infty\).
\end{theorem}

\begin{theorem}[Main theorem]
\label{ch05-thm-main}
\source{chow-knopf-ricci-flow:page-0118, chow-knopf-ricci-flow:page-0119}
\dcref{Theorem 5.1}
\group{main-theorem}

If \((\mathcal{M}^2,g_0)\) is a closed Riemannian surface, there exists a unique solution \(g(t)\) of the normalized Ricci flow
\[
\frac{\partial}{\partial t}g=(r-R)g,\qquad g(0)=g_0,
\]
which exists for all time and converges uniformly in any \(C^k\)-norm to a smooth metric \(g_\infty\) of constant curvature.
\end{theorem}
